% ================================================================== Chapitre 6
\chapter{Influence des paramètres et optimisation}
\label{ch:vanthoff}

Une fois l'équilibre établi, on peut le \emph{déplacer} en jouant sur les conditions opératoires. Deux leviers existent, et toute la discussion s'organise autour d'eux : modifier $K^\circ$ (seule la température le permet) ou modifier $Q_r$ (pression, composition).

\section{Influence de la température : relation de Van~'t~Hoff}

\begin{resultat}{Relation de Van~'t~Hoff}
\[ \boxed{\frac{\dd\ln K^\circ}{\dd T} = \frac{\Delta_rH^\circ}{RT^2}} \]
\begin{itemize}
  \item réaction \textbf{endothermique} ($\Delta_rH^\circ > 0$) : $K^\circ$ \emph{augmente} avec $T$ ;
  \item réaction \textbf{exothermique} ($\Delta_rH^\circ < 0$) : $K^\circ$ \emph{diminue} quand $T$ augmente ;
  \item réaction athermique : $K^\circ$ quasi indépendante de $T$.
\end{itemize}
\end{resultat}

\begin{demo}[Démonstration dans l'approximation d'Ellingham]
$\ln K^\circ = -\dfrac{\Delta_rG^\circ}{RT} = -\dfrac{\Delta_rH^\circ}{RT} + \dfrac{\Delta_rS^\circ}{R}$. Avec $\Delta_rH^\circ$ et $\Delta_rS^\circ$ constantes, la dérivée par rapport à $T$ vaut $\dfrac{\Delta_rH^\circ}{RT^2}$.
\end{demo}

\begin{horsprog}[validité générale]
La relation de Van~'t~Hoff reste vraie même si $\Delta_rH^\circ$ dépend de $T$. On le montre avec la relation de Gibbs--Helmholtz $\dfrac{\dd}{\dd T}\left(\dfrac{\Delta_rG^\circ}{T}\right) = -\dfrac{\Delta_rH^\circ}{T^2}$, conséquence de $\left(\partial G/\partial T\right)_P = -S$.
\end{horsprog}

\begin{resultat}{Forme intégrée ($\Delta_rH^\circ$ constante)}
\[ \ln\frac{K^\circ(T_2)}{K^\circ(T_1)} = -\frac{\Delta_rH^\circ}{R}\left(\frac{1}{T_2}-\frac{1}{T_1}\right) \]
$\ln K^\circ$ est une fonction affine de $1/T$, de pente $-\Delta_rH^\circ/R$.
\end{resultat}

\begin{figure}[H]
\centering
\begin{tikzpicture}
\begin{axis}[cours,width=0.74\linewidth,height=0.42\linewidth,xlabel={$1000/T$ ($\mathrm{K^{-1}}$)},ylabel={$\ln K^\circ$},
  xmin=1.0,xmax=2.6,ymin=-12,ymax=3,clip=false]
  \addplot[c1,mark=*,mark size=1.6pt] table[x=invT_milli,y=lnK]{data/nh3_vh.dat};
  \draw[c3,line width=0.8pt] (axis cs:1.6,-6.16) -- (axis cs:2.0,-6.16) -- (axis cs:2.0,-1.74);
  \node[c3,font=\small,anchor=west,align=left] at (axis cs:2.04,-5.6) {pente $= -\Delta_rH^\circ/R$\\$= +\SI{1,10e4}{K}$};
  \node[font=\small,anchor=south east,align=right] at (axis cs:2.58,-11.6) {\ce{N2 + 3 H2 = 2 NH3}, exothermique :\\$K^\circ$ diminue quand $T$ augmente};
\end{axis}
\end{tikzpicture}
\caption{Droite de Van~'t~Hoff pour la synthèse de l'ammoniac (approximation d'Ellingham), de $\SI{400}{K}$ à $\SI{900}{K}$.}
\label{fig:vh}
\end{figure}

\begin{methode}[calculer $K^\circ$ à une température quelconque]
\begin{itemize}
  \item \emph{Voie 1} : calculer $\Delta_rG^\circ(T) = \Delta_rH^\circ - T\Delta_rS^\circ$ avec les tables, puis $K^\circ(T) = \exp(-\Delta_rG^\circ/RT)$.
  \item \emph{Voie 2} : partir d'une valeur connue $K^\circ(T_1)$ et utiliser la forme intégrée de Van~'t~Hoff.
\end{itemize}
Les deux voies donnent le même résultat dans l'approximation d'Ellingham.
\end{methode}

\begin{exemple}[ammoniac à $\SI{450}{\celsius}$]
Voie 1 : $\Delta_rG^\circ(\SI{723}{K}) = -91{,}8 + 723\times0{,}1981 = \SI{51,4}{kJ.mol^{-1}}$, d'où $K^\circ = \mathrm{e}^{-51{,}4\times10^3/(8{,}314\times723)} = \num{1,9e-4}$.

\noindent Voie 2 : $K^\circ(723) = \num{5,4e5}\times\exp\left[\dfrac{91{,}8\times10^3}{8{,}314}\left(\dfrac{1}{723}-\dfrac{1}{298}\right)\right] = \num{1,9e-4}$.

\noindent En passant de $\SI{25}{\celsius}$ à $\SI{450}{\celsius}$, $K^\circ$ est divisée par environ $3\times10^9$.
\end{exemple}

\begin{methode}[mesurer une enthalpie standard de réaction]
On mesure $K^\circ$ à plusieurs températures (par titrage, conductimétrie, mesure de pression\dots), on trace $\ln K^\circ$ en fonction de $1/T$ et on fait une régression linéaire : la pente vaut $-\Delta_rH^\circ/R$ et l'ordonnée à l'origine $\Delta_rS^\circ/R$. C'est aussi une manière de vérifier l'approximation d'Ellingham (les points sont-ils alignés ?).
\end{methode}

\section{Optimisation par modification de $K^\circ$ : choix de la température}

Seule la température modifie $K^\circ$. Pour \emph{augmenter le rendement} à l'équilibre :
\begin{itemize}
  \item réaction exothermique : travailler à \emph{basse} température ;
  \item réaction endothermique : travailler à \emph{haute} température.
\end{itemize}
Mais la vitesse de réaction augmente toujours avec $T$ (loi d'Arrhenius). Pour une réaction exothermique, il y a donc un \textbf{compromis} entre thermodynamique (basse $T$) et cinétique (haute $T$), que l'on améliore avec un catalyseur (qui accélère la réaction sans modifier $K^\circ$).

\section{Optimisation par modification de $Q_r$}

On part d'un système à l'équilibre ($Q_r = K^\circ$). On le perturbe à $T$ fixée, \emph{instantanément}, puis on regarde comment $Q_r$ a varié : le système évolue ensuite de façon à ramener $Q_r$ vers $K^\circ$.

\begin{resultat}{Quotient de réaction d'un mélange gazeux}
Pour des gaz parfaits, $a_i = \dfrac{P_i}{P^\circ} = \dfrac{n_i}{n_{\text{tot,gaz}}}\,\dfrac{P}{P^\circ}$, d'où, en notant $\Delta\nu_{\text{gaz}} = \sum_{\text{gaz}}\nu_i$ :
\[ Q_r = \Big(\prod_{\text{gaz}}n_i^{\,\nu_i}\Big)\left(\frac{P}{n_{\text{tot,gaz}}\,P^\circ}\right)^{\Delta\nu_{\text{gaz}}}\times\big(\text{activités des autres constituants}\big). \]
\end{resultat}

\begin{methode}[prévoir l'effet d'une perturbation]
\begin{enumerate}
  \item Exprimer $Q_r$ en fonction de la grandeur modifiée ($P$, $n_i$, $n_{\text{tot}}$\dots), les autres restant fixes.
  \item Dire si la perturbation fait augmenter ou diminuer $Q_r$ à partir de $Q_r = K^\circ$.
  \item Conclure : $Q_r < K^\circ$ entraîne une évolution dans le sens direct ; $Q_r > K^\circ$ dans le sens indirect.
\end{enumerate}
\end{methode}

\begin{resultat}{Effet des principaux paramètres (gaz parfaits, $T$ fixée)}
\renewcommand{\arraystretch}{1.35}
\begin{tabularx}{\linewidth}{>{\raggedright\arraybackslash}p{4.4cm} X}
\toprule
\textbf{Perturbation} & \textbf{Effet sur $Q_r$ et conséquence} \\ \midrule
augmentation de la pression totale $P$ & $Q_r\propto P^{\Delta\nu_{\text{gaz}}}$ : si $\Delta\nu_{\text{gaz}}<0$, $Q_r$ diminue $\Rightarrow$ sens direct (vers \emph{moins} de gaz) ; si $\Delta\nu_{\text{gaz}} = 0$, aucun effet \\
ajout d'un gaz inerte à $P$ fixée & $n_{\text{tot,gaz}}$ augmente : $Q_r\propto n_{\text{tot,gaz}}^{-\Delta\nu_{\text{gaz}}}$. Si $\Delta\nu_{\text{gaz}}<0$, $Q_r$ augmente $\Rightarrow$ sens indirect (même effet qu'une baisse de $P$) \\
ajout d'un gaz inerte à $V$ fixé & les pressions partielles $P_i = n_iRT/V$ sont inchangées : aucun effet \\
élimination d'un produit (distillation, liquéfaction, précipitation) & $Q_r$ diminue $\Rightarrow$ sens direct \\
ajout d'un réactif (en solution, ou gaz à $V$ fixé) & $Q_r$ diminue $\Rightarrow$ sens direct \\
ajout ou retrait d'un solide ou liquide pur (sans le faire disparaître) & activité égale à 1, $Q_r$ inchangé : aucun effet \\
\bottomrule
\end{tabularx}
\end{resultat}

\begin{pieges}
\begin{itemize}
  \item Raisonner sur $Q_r$, pas sur des intuitions : l'ajout d'un gaz inerte n'a pas le même effet à $P$ fixée et à $V$ fixé.
  \item Un catalyseur ne modifie ni $K^\circ$ ni $Q_r$ : il ne change pas l'état d'équilibre, il permet seulement de l'atteindre plus vite.
\end{itemize}
\end{pieges}

\begin{horsprog}[principe de modération de Le Chatelier]
Les résultats du tableau se résument qualitativement ainsi : « le système réagit à une perturbation en s'opposant à celle-ci » (une hausse de $P$ favorise le sens qui consomme du gaz, une hausse de $T$ favorise le sens endothermique). C'est un moyen mnémotechnique ; aux concours, on attend la justification par $Q_r$ et $K^\circ$.
\end{horsprog}

\section{Application : la synthèse industrielle de l'ammoniac}

Le procédé Haber--Bosch réalise \ce{N2(g) + 3 H2(g) = 2 NH3(g)}, avec $\Delta_rH^\circ = \SI{-91,8}{kJ.mol^{-1}}$ (exothermique) et $\Delta\nu_{\text{gaz}} = -2$.

\begin{itemize}
  \item \textbf{Température} : la thermodynamique demande une basse température, mais la réaction est alors trop lente. Compromis : 400 à $\SI{500}{\celsius}$ avec un catalyseur à base de fer.
  \item \textbf{Pression} : $\Delta\nu_{\text{gaz}}<0$, donc une haute pression favorise la synthèse. On travaille vers 150 à $\SI{300}{bar}$ (au-delà, le coût des installations devient prohibitif).
  \item \textbf{Composition} : on élimine \ce{NH3} par liquéfaction à la sortie du réacteur, et l'on recycle \ce{N2} et \ce{H2} non transformés. On évite l'accumulation de gaz inertes (argon de l'air), qui défavorise la synthèse à $P$ fixée.
\end{itemize}

\begin{exemple}[effet de la pression à $\SI{723}{K}$]
Mélange initial stœchiométrique (\SI{1}{mol} de \ce{N2}, \SI{3}{mol} de \ce{H2}) ; on résout $Q_r(\xi) = K^\circ = \num{1,9e-4}$ avec $n_{\text{tot}} = 4 - 2\xi$ :
\[ \frac{(2\xi)^2(4-2\xi)^2}{(1-\xi)(3-3\xi)^3}\left(\frac{P^\circ}{P}\right)^2 = K^\circ \]
\begin{center}
\renewcommand{\arraystretch}{1.25}
\begin{tabular}{lcccc}
\toprule
$P$ (bar) & 1 & 50 & 200 & 300 \\ \midrule
$\xi_{\text{éq}}$ (mol) & 0,009 & 0,27 & 0,53 & 0,60 \\
fraction molaire de \ce{NH3} & 0,004 & 0,16 & 0,36 & 0,43 \\
\bottomrule
\end{tabular}
\end{center}
Sous $\SI{1}{bar}$ à $\SI{450}{\celsius}$, la synthèse est quasi nulle ; à $\SI{200}{bar}$, plus du tiers du mélange est de l'ammoniac. (Valeurs calculées numériquement par résolution de l'équation ci-dessus.)
\end{exemple}
